\section*{按顶刊/顶会标准的逐节修改清单（实验增强重点版）}

\subsection*{总体实验判断}
\begin{itemize}
    \item 当前稿件在写作层面已经相当成熟，但实验层面的核心问题仍然是：\textbf{现有实验矩阵足以证明“当前不能宣称 gain”，却还不足以回答“为什么不能”以及“下一步最小可决定实验是什么”。}
    \item 结合主文与附录，当前实验已经明确暴露出三项事实：
    \begin{enumerate}
        \item paired validation split 仅有 4 张图像，导致结果对单张样本极其敏感；
        \item late fusion、same-modality KD、cross-modality KD 在 seed-aggregated confusion 上都表现出将全部验证样本预测为 positive 的倾向；
        \item 当前 MHRA 未被验证，且在 module ablation 中，去掉 MHRA 反而得到略高均值。
    \end{enumerate}
    \item 因此，实验部分的下一步不应再以“更多微型 ablation”为主，而应转向 \textbf{瓶颈定位实验（bottleneck-isolation experiments）}、\textbf{最小决定性实验（minimum decisive experiments）} 和 \textbf{证据门槛实验（evidence-threshold experiments）}。
\end{itemize}

\subsection*{1. Abstract}
\begin{itemize}
    \item \textbf{实验相关优点}
    \begin{itemize}
        \item 摘要已经正确说明：student-only baseline 最强，late fusion 与 cross-modality KD 未显示稳定增益，MHRA 也未被验证。
    \end{itemize}
    \item \textbf{实验相关仍需增强的点}
    \begin{itemize}
        \item 当前摘要仍未点出“\textbf{缺失的关键实验类型}”。
        \item 顶刊风格下，摘要不一定要列出未来所有实验，但应明确让读者知道：当前结论之所以是 evidence-bounded，是因为\textbf{最小 cross-modal baseline 尚未被充分隔离验证，paired cohort 也尚未达到可解释规模。}
    \end{itemize}
    \item \textbf{建议}
    \begin{itemize}
        \item 在摘要结尾加入一句关于实验边界的话，例如：
        \begin{quote}
            当前结果的解释力受限于极小 paired validation support 和缺少 stripped-down cross-modal control 的实验边界，因此本文更接近 benchmark scaffold 而非最终方法结论。
        \end{quote}
    \end{itemize}
\end{itemize}

\subsection*{2. Introduction}
\begin{itemize}
    \item \textbf{实验相关仍需增强的点}
    \begin{itemize}
        \item 引言已说明该研究设定有意义，但还可以更进一步指出：\textbf{本问题之所以长期容易被误判，是因为缺少一套标准化实验矩阵来区分 data scarcity、modality gap、baseline insufficiency 和 module over-parameterization。}
    \end{itemize}
    \item \textbf{建议}
    \begin{itemize}
        \item 在引言末尾增加一句，把实验目标定义为：
        \begin{quote}
            本文不仅评估一个 cross-modality distillation 实例，也尝试定义一个最小实验框架，用于判断 observed gains 到底来自 cross-modality supervision、architecture complexity，还是仅仅来自样本偏差与不稳定验证支持。
        \end{quote}
    \end{itemize}
\end{itemize}

\subsection*{3. Methodology}
\begin{itemize}
    \item \textbf{实验相关仍需增强的点}
    \begin{itemize}
        \item 当前方法部分已经通过 hypothesis-driven 叙述减轻了“模块堆叠”问题，但还可以更进一步为实验设计服务。
        \item 具体来说，Method 应该更明确地告诉读者：\textbf{哪些模块对应哪些可检验假设，哪些实验是为验证这些假设而设计的。}
    \end{itemize}
    \item \textbf{建议}
    \begin{itemize}
        \item 在 3.3 或 3.4 后加入一个小段落，命名为：
        \begin{quote}
            \textbf{Testable hypotheses and corresponding controls}
        \end{quote}
        \item 明确列出：
        \begin{enumerate}
            \item H1: CT supervision can help X-ray deployment beyond student-only training；
            \item H2: the benefit, if any, is not reducible to same-modality KD alone；
            \item H3: the benefit, if any, is not reducible to inference-time late fusion；
            \item H4: DPE/MHRA/DFPN improve transfer beyond stripped-down cross-modal logit KD；
            \item H5: observed gains are not artifacts of class imbalance or tiny validation support.
        \end{enumerate}
        \item 这样后续实验节就能围绕假设验证展开，而不是看起来像若干散乱对比。
    \end{itemize}
\end{itemize}

\subsection*{4. Experiments: 当前最缺的实验类型}
\begin{itemize}
    \item \textbf{当前最缺失的不是更多超参数 ablation，而是以下五类实验：}
    \begin{enumerate}
        \item \textbf{最小跨模态基线实验（Minimal cross-modal baseline experiment）}
        \item \textbf{同病例五路对照实验（Same-case five-way comparison）}
        \item \textbf{数据规模敏感性实验（Paired-cohort scaling experiment）}
        \item \textbf{类别不平衡/阈值稳定性实验（Imbalance and threshold-behavior experiment）}
        \item \textbf{模块必要性递进实验（Progressive module-necessity experiment）}
    \end{enumerate}
\end{itemize}

\subsection*{4.1 最小跨模态基线实验（当前最重要）}
\begin{itemize}
    \item \textbf{为什么缺}
    \begin{itemize}
        \item 当前主文已承认：缺少一个 stripped-down cross-modal logit-KD baseline，而现有 student-only / same-modality KD / ablation 只能构成“近似替代”，无法完全回答“到底是 cross-modality supervision 本身无效，还是当前模块栈掩盖了它的可能收益”。:contentReference[oaicite:2]{index=2}
    \end{itemize}

    \item \textbf{应该怎么设计}
    \begin{itemize}
        \item 保留完全相同的数据划分、teacher/student backbone、训练轮数、优化器、temperature 和 $\alpha$；
        \item 去掉 DPE、MHRA、DFPN，仅保留：
        \begin{enumerate}
            \item CT teacher logits；
            \item X-ray student logits；
            \item hard CE + soft KL。
        \end{enumerate}
        \item 即构造一个真正的 \textbf{plain cross-modal logit KD} baseline。
    \end{itemize}

    \item \textbf{应该报告什么}
    \begin{itemize}
        \item 与以下四个系统在完全相同病例上做 repeated-run 对比：
        \begin{enumerate}
            \item student-only X-ray；
            \item same-modality KD；
            \item plain cross-modal logit KD；
            \item current JDCNet；
            \item late fusion。
        \end{enumerate}
        \item 指标至少包括：Accuracy、Macro-F1、Balanced Accuracy、Specificity、Recall。
        \item 不建议再把 ROC-AUC 作为 headline 指标，因为在当前四图验证集上它解释性很弱。
    \end{itemize}

    \item \textbf{如何实施}
    \begin{itemize}
        \item 在当前仓库中新增一个最小配置开关，例如：
        \begin{quote}
            \texttt{--crossmodal\_mode plain\_logit\_kd}
        \end{quote}
        \item 保证除了模块关闭外，其它训练条件不变。
        \item 所有实验必须与现有 repeated-run pipeline 共用同一套 manifest 与 seed，避免隐性数据漂移。
    \end{itemize}

    \item \textbf{这个实验回答什么}
    \begin{itemize}
        \item 若 plain cross-modal logit KD 仍不优于 student-only，则说明当前瓶颈大概率在 \textbf{data scarcity / modality gap}；
        \item 若 plain cross-modal logit KD 优于 current JDCNet，则说明当前问题更可能在 \textbf{module over-parameterization or unstable feature manipulation}；
        \item 若 plain cross-modal logit KD 与 same-modality KD 接近，则说明跨模态教师信号未明显超出 teacher regularization 本身。
    \end{itemize}
\end{itemize}

\subsection*{4.2 同病例五路对照实验（Minimum decisive experiment）}
\begin{itemize}
    \item \textbf{为什么缺}
    \begin{itemize}
        \item 你现在虽然在主文与附录中已经提出了 five-way comparison 的概念，但它还没有真正落地为一节完整实验设计。
    \end{itemize}

    \item \textbf{应该怎么设计}
    \begin{itemize}
        \item 使用同一 patient-level paired cohort、同一训练/验证病例；
        \item 五个系统固定为：
        \begin{enumerate}
            \item student-only X-ray；
            \item inference-time late fusion；
            \item same-modality KD；
            \item stripped-down cross-modal logit KD；
            \item current JDCNet。
        \end{enumerate}
        \item 每个系统都做完全相同的 repeated-run seeds。
    \end{itemize}

    \item \textbf{应该报告什么}
    \begin{itemize}
        \item 除均值和标准差外，还应报告：
        \begin{itemize}
            \item per-seed case-level prediction table；
            \item seed-aggregated confusion summaries；
            \item 若样本量增大，可增加 bootstrap CI。
        \end{itemize}
        \item 图表上建议保留你现在的 repeated-run bar + per-seed points 风格，因为它很适合展示 instability。
    \end{itemize}

    \item \textbf{这个实验回答什么}
    \begin{itemize}
        \item 它能最直接地区分：
        \begin{itemize}
            \item CT supervision 本身是否有价值；
            \item 价值是否只在 inference-time fusion 下出现；
            \item 模块栈是否真的优于 stripped-down baseline。
        \end{itemize}
    \end{itemize}
\end{itemize}

\subsection*{4.3 数据规模敏感性实验（Paired-cohort scaling experiment）}
\begin{itemize}
    \item \textbf{为什么缺}
    \begin{itemize}
        \item 目前论文已经证明 tiny paired cohort 是瓶颈，但还没有实验量化“\textbf{到底多大才开始出现稳定性改善}”。
    \end{itemize}

    \item \textbf{应该怎么设计}
    \begin{itemize}
        \item 如果数据允许，把 paired cohort 按 patient 数量做多档构造，例如：
        \begin{enumerate}
            \item 25\% paired patients；
            \item 50\% paired patients；
            \item 75\% paired patients；
            \item 100\% paired patients。
        \end{enumerate}
        \item 若原始 paired cases 不够，可以换成更宽松的 temporal matching，但必须单独报告匹配规则变化。
    \end{itemize}

    \item \textbf{应该报告什么}
    \begin{itemize}
        \item 各档 paired size 下五个系统的 mean/std；
        \item instability 曲线，例如 Macro-F1 std vs paired size；
        \item Specificity 与 Balanced Accuracy 随 paired size 的变化。
    \end{itemize}

    \item \textbf{这个实验回答什么}
    \begin{itemize}
        \item 它能回答：
        \begin{quote}
            当前 failure 到底是结构性 modality gap，还是因为 paired support 还没有达到最小有效规模。
        \end{quote}
        \item 这是顶刊评审非常在意的问题，因为它决定这项研究是“方向不成立”，还是“方向成立但证据未到位”。
    \end{itemize}
\end{itemize}

\subsection*{4.4 类别不平衡与阈值稳定性实验}
\begin{itemize}
    \item \textbf{为什么缺}
    \begin{itemize}
        \item 当前附录已经非常强地表明存在 prevalence-biased positive prediction：late fusion、same-modality KD、cross-modality KD 在 seed-aggregated confusion 下都把验证图像全部判为 positive。:contentReference[oaicite:4]{index=4}
        \item 但还没有实验性地回答：这种现象是来自训练集不平衡、logit calibration 失衡、还是 teacher signal 本身偏置。
    \end{itemize}

    \item \textbf{应该怎么设计}
    \begin{itemize}
        \item 至少增加两类 control：
        \begin{enumerate}
            \item \textbf{class-balanced training control}：对训练集使用 class-balanced sampler 或 loss weighting；
            \item \textbf{threshold calibration control}：验证集上不固定 0.5 阈值，而报告基于训练折或小验证折校准后的阈值表现。
        \end{enumerate}
        \item 若担心 validation 太小，可用 training split 内部再做 calibration fold，不能直接用 test-like validation 反复调阈值。
    \end{itemize}

    \item \textbf{应该报告什么}
    \begin{itemize}
        \item Specificity、Recall、Balanced Accuracy 必须成为这一实验的主指标；
        \item confusion matrix 必须与默认设置并排报告；
        \item 最好增加 prediction score histogram 或 positive-rate statistics。
    \end{itemize}

    \item \textbf{这个实验回答什么}
    \begin{itemize}
        \item 如果 class balancing/calibration 后 cross-modal KD 仍无优势，则说明问题更接近 \textbf{true transfer failure}；
        \item 如果 balancing/calibration 后才出现 improvement，则说明此前的 negative result 很大程度来自 \textbf{imbalance-dominated decision behavior}。
    \end{itemize}
\end{itemize}

\subsection*{4.5 模块必要性递进实验（比当前 ablation 更强）}
\begin{itemize}
    \item \textbf{为什么缺}
    \begin{itemize}
        \item 你现在已经有 w/o DPE、w/o MHRA、w/o DFPN 的 ablation，但这类“单模块移除”还不够回答模块栈到底哪里开始变得不必要。
    \end{itemize}

    \item \textbf{应该怎么设计}
    \begin{itemize}
        \item 按递进复杂度做实验，而不是只做 leave-one-out：
        \begin{enumerate}
            \item plain cross-modal logit KD；
            \item + DPE；
            \item + DPE + DFPN；
            \item + DPE + MHRA；
            \item + DPE + MHRA + DFPN（即 current JDCNet）。
        \end{enumerate}
        \item 这比当前 ablation 更能回答“复杂度是何时开始引入的、是否值得”。
    \end{itemize}

    \item \textbf{应该报告什么}
    \begin{itemize}
        \item 除 mean/std 外，建议报告参数量、训练时间、每阶段性能变化；
        \item 若可能，增加 one-paragraph analysis：复杂度增加是否换来任何稳定性改善。
    \end{itemize}

    \item \textbf{这个实验回答什么}
    \begin{itemize}
        \item 它能把当前“MHRA 似乎没用”提升为更系统的结论：
        \begin{quote}
            到底是某个模块无效，还是整个 feature-level complexity 在当前证据条件下都不值得。
        \end{quote}
    \end{itemize}
\end{itemize}

\subsection*{4.6 外部可迁移性实验（可选，但顶刊很加分）}
\begin{itemize}
    \item \textbf{为什么缺}
    \begin{itemize}
        \item 当前主文已说明 Kaggle 下载的数据不是 same-patient chest X-ray/CT benchmark，因此没纳入主结果，这是合理的。:contentReference[oaicite:5]{index=5}
        \item 但从顶刊角度，如果能做一个弱形式 external sanity check，会更强。
    \end{itemize}

    \item \textbf{应该怎么设计}
    \begin{itemize}
        \item 不把 external cohort 当 paired benchmark；
        \item 仅把在 paired cohort 上训练好的 student-only 与 cross-modal variant 拿去外部 X-ray 数据上做 out-of-cohort inference sanity check；
        \item 不做强 claim，只看两者是否有一致退化。
    \end{itemize}

    \item \textbf{应该报告什么}
    \begin{itemize}
        \item 只需报告 Accuracy、Macro-F1、Balanced Accuracy；
        \item 必须明确写明：这不是配对跨模态验证，只是迁移稳定性 sanity check。
    \end{itemize}

    \item \textbf{这个实验回答什么}
    \begin{itemize}
        \item 它能让你更有底气说明：当前负结果并非完全局限于单一 split artifact。
    \end{itemize}
\end{itemize}

\subsection*{5. Results 部分应新增的实验叙事}
\begin{itemize}
    \item 结果部分不应只写“跑了哪些实验”，而应明确按以下层次组织：
    \begin{enumerate}
        \item \textbf{Can CT supervision help at all?} $\rightarrow$ student-only vs plain cross-modal KD；
        \item \textbf{Is any benefit modality-specific or just teacher regularization?} $\rightarrow$ same-modality KD vs cross-modal KD；
        \item \textbf{Is complexity justified?} $\rightarrow$ plain cross-modal KD vs current JDCNet vs progressive module addition；
        \item \textbf{Are observed behaviors dominated by data imbalance?} $\rightarrow$ class-balanced / calibrated controls；
        \item \textbf{Does the conclusion change as paired evidence increases?} $\rightarrow$ scaling experiment。
    \end{enumerate}
    \item 只有这样，实验节才不只是“结果表集合”，而是真正回答科学问题。
\end{itemize}

\subsection*{6. 实施层面建议（非常具体）}
\begin{itemize}
    \item \textbf{实验执行顺序建议}
    \begin{enumerate}
        \item 先做 \textbf{plain cross-modal logit KD baseline}；
        \item 再做 \textbf{same-case five-way comparison}；
        \item 再做 \textbf{class-balanced / threshold-calibrated controls}；
        \item 再做 \textbf{progressive module-necessity experiment}；
        \item 最后做 \textbf{paired-cohort scaling experiment}。
    \end{enumerate}

    \item \textbf{优先级理由}
    \begin{itemize}
        \item 第 1 和第 2 项是最小决定性实验；
        \item 第 3 项回答 prevalence bias；
        \item 第 4 项回答 complexity 是否值得；
        \item 第 5 项回答数据规模是否是主瓶颈。
    \end{itemize}

    \item \textbf{代码实施建议}
    \begin{itemize}
        \item 所有新实验必须复用现有 manifest generation 和 repeated-run pipeline；
        \item 所有新配置建议通过 config flag 实现，不要复制一套新训练脚本；
        \item 每新增一种实验模式，都输出：
        \begin{itemize}
            \item per-seed csv；
            \item repeated-run summary csv；
            \item confusion summary；
            \item figure export。
        \end{itemize}
        \item 避免手工后处理，保证可复验。
    \end{itemize}
\end{itemize}

\subsection*{7. 按优先级排序的实验补强任务}
\begin{enumerate}
    \item \textbf{最高优先级：}补上 stripped-down plain cross-modal logit-KD baseline。
    \item \textbf{最高优先级：}在同一 paired cases 上完成 five-way repeated-run comparison。
    \item \textbf{高优先级：}加入 imbalance-control 和 threshold-behavior experiment。
    \item \textbf{高优先级：}把现有 leave-one-out ablation 升级成 progressive module-complexity experiment。
    \item \textbf{中优先级：}做 paired-cohort scaling experiment。
    \item \textbf{可选增强：}做 external X-ray sanity-check experiment。
\end{enumerate}

\subsection*{8. 最终实验判断}
\begin{itemize}
    \item 当前稿件已经非常成功地证明了：\textbf{在现有 tiny paired cohort 下，强 claim 不成立。}
    \item 但若要从“谨慎负结果论文”进一步升级到“顶刊/顶会更强说服力论文”，最关键的不是再润色文字，而是补上：
    \begin{quote}
        能够隔离 cross-modality supervision 本身、teacher regularization、本征类别偏差、以及模块复杂度四者关系的最小决定性实验矩阵。
    \end{quote}
    \item 一旦这些实验补齐，你的论文就不仅能说“当前不能宣称 gain”，还可以更强地说：
    \begin{quote}
        我们已经知道 gain 为什么还不能被宣称，以及未来什么样的实验才有资格支持它。
    \end{quote}
\end{itemize}